\documentclass{article}
\usepackage[T1]{fontenc}
\usepackage{lmodern}
\usepackage{graphicx}
\usepackage{amsmath,amssymb}
\usepackage[a4paper,margin=2cm]{geometry}
\usepackage[absolute,overlay]{textpos}
\setlength{\TPHorizModule}{1mm}
\setlength{\TPVertModule}{1mm}
\usepackage[indonesian]{babel}
\usepackage{hyperref}
\setlength{\emergencystretch}{2em}
% unit: Halaman 25

\begin{document}

% === Halaman 25 (llm) ===

\section*{Cipher Alir (Stream Cipher)}

\begin{itemize}
    \item Mengenkripsi plainteks menjadi chiperteks dengan cara:
    \begin{itemize}
        \item $1 \text{ bit plainteks} \oplus 1 \text{ bit kunci}$, atau
        \item $1 \text{ byte plainteks} \oplus 1 \text{ byte skunci}$, atau
        \item $n \text{ bit plainteks} \oplus n \text{ bit kunci}$
    \end{itemize}
    setiap kali enkripsi
    
    \item Dekripsi dilakukan dengan cara:
    \begin{itemize}
        \item $1 \text{ bit cipherteks} \oplus 1 \text{ bit kunci}$, atau
        \item $1 \text{ byte cipherteks} \oplus 1 \text{ byte skunci}$, atau
        \item $n \text{ bit cipherteks} \oplus n \text{ bit kunci}$
    \end{itemize}
    setiap kali dekripsi.
\end{itemize}

\end{document}
